\documentclass[11pt]{article}

\usepackage[margin=1in]{geometry}
\usepackage{amsmath,amssymb,amsthm,mathtools}
\usepackage{enumitem}
\usepackage{hyperref}
\usepackage{xcolor}
\usepackage{microtype}

\hypersetup{
  colorlinks=true,
  linkcolor=blue!60!black,
  citecolor=blue!60!black,
  urlcolor=blue!60!black
}

\newtheorem{theorem}{Theorem}[section]
\newtheorem{proposition}[theorem]{Proposition}
\newtheorem{lemma}[theorem]{Lemma}
\newtheorem{definition}[theorem]{Definition}
\newtheorem{remark}[theorem]{Remark}

\newcommand{\Conf}{\operatorname{Conf}}
\newcommand{\dR}{\operatorname{dR}}
\newcommand{\Span}{\operatorname{Span}}

\title{A Lean/SymPy Formal Skeleton for Three-Point Non-Isotropic Configurations over a Complex Quadric}
\author{Nick Goutev \and Lean/SymPy formalization audit}
\date{June 2026}

\begin{document}
\maketitle

\begin{abstract}
Let $V=\mathbb C^D$ with $D$ even, and let $q$ be a nondegenerate complex quadratic form.  We consider the configuration space of three ordered points whose pairwise separations are non-isotropic:
\[
F_3^{\mathrm{noniso}}(V,q)=
\{(x_1,x_2,x_3)\in V^3: q(x_i-x_j)\neq 0\text{ for }i<j\}.
\]
This note records a theorem-honest finite formalization in Lean~4 and a supporting SymPy audit.  The Lean development proves the finite combinatorial target needed for the de Rham computation: translation-reduced coordinates, formal edge generators $\alpha_{ij}$ and $\beta_{ij}$ with degrees $1$ and $D-1$, the arity-three cooperad edge bookkeeping, and a degree-one Orlik--Solomon skeleton with the Arnold relation
\[
A_{12}A_{23}-A_{12}A_{13}+A_{23}A_{13}=0.
\]
The analytic comparison identifying this formal skeleton with the actual de Rham cohomology of the quadric-divisor complement is deliberately left as an explicit socket.
\end{abstract}

\section{Configuration Space}
Let
\[
V=\mathbb C^D,
\qquad
q(x)=\sum_{a=1}^D x_a^2.
\]
Over $\mathbb C$, every nondegenerate quadratic form is linearly equivalent to this standard model.  The three-point non-isotropic configuration space is
\[
F_3^{\mathrm{noniso}}(V,q)=
\left\{(x_1,x_2,x_3)\in V^3:
q(x_1-x_2)q(x_1-x_3)q(x_2-x_3)\neq 0
\right\}.
\]
The expected de Rham calculation is governed by the complement of the three quadratic divisors
\[
q(x_1-x_2)=0,
\qquad
q(x_1-x_3)=0,
\qquad
q(x_2-x_3)=0.
\]

\section{Translation Reduction}
Choose $x_1$ as basepoint and write
\[
u=x_1-x_2,
\qquad
v=x_1-x_3.
\]
Then
\[
x_2=x_1-u,
\qquad
x_3=x_1-v,
\]
and the three non-isotropic conditions become
\[
q(u)\neq 0,
\qquad
q(v)\neq 0,
\qquad
q(v-u)\neq 0.
\]
Thus the configuration space splits, at the level of coordinates, into a translation factor and a reduced two-vector complement:
\[
F_3^{\mathrm{noniso}}(V,q)
\simeq
V\times
\{(u,v)\in V^2:q(u)q(v)q(v-u)\neq 0\}.
\]

\begin{proposition}[Lean formalized translation bookkeeping]
The Lean file \texttt{proofs/QuadraticConfiguration3.lean} defines structures
\[
\texttt{Config3}(D),
\qquad
\texttt{Reduced3}(D),
\]
with maps
\[
\texttt{Config3.toReduced},
\qquad
\texttt{Reduced3.toConfig},
\]
which preserve the three non-isotropic predicates.
\end{proposition}

\section{Formal Edge Generators}
Let $K_3$ be the complete graph on vertices $\{1,2,3\}$, with edges
\[
12,
\qquad
13,
\qquad
23.
\]
The formal cohomology target uses two generator types on each edge:
\[
\alpha_{ij},
\qquad
\beta_{ij}.
\]
Their degrees are
\[
|\alpha_{ij}|=1,
\qquad
|\beta_{ij}|=D-1.
\]
The $\alpha$ classes correspond to logarithmic classes $d\log q(x_i-x_j)$, while the $\beta$ classes denote the additional quadric angular class expected in even complex dimension.  In the current Lean formalization, this is a formal generator/degrees layer; the analytic construction of $\beta$ is socketed.

\begin{theorem}[Lean formalized generator degrees]
For every edge $ij\in\{12,13,23\}$, the formal degree function satisfies
\[
\deg(\alpha_{ij})=1,
\qquad
\deg(\beta_{ij})=D-1.
\]
\end{theorem}

\section{Arity-Three Cooperad Bookkeeping}
The arity-three cooperad has three elementary decompositions:
\[
(12)|3,
\qquad
(13)|2,
\qquad
(23)|1.
\]
For such a decomposition, the unique edge internal to the two-point cluster maps to the internal factor.  The two remaining edges cross from the cluster to the outside vertex and map to the outer factor.

For example, for $(23)|1$,
\[
23\longmapsto \mathrm{internal},
\qquad
12\longmapsto \mathrm{outer},
\qquad
13\longmapsto \mathrm{outer}.
\]
The generator label is preserved:
\[
\alpha\mapsto\alpha,
\qquad
\beta\mapsto\beta.
\]

\begin{theorem}[Lean formalized arity-three cooperad edge rule]
For each block decomposition of type $(ij)|k$, there exists exactly one internal edge.  Every edge maps either to the internal factor or the outer factor, and the generator kind $\alpha/\beta$ is preserved.
\end{theorem}

\section{Orlik--Solomon Skeleton}
The degree-one Orlik--Solomon-style skeleton uses three formal generators
\[
A_{12},
\qquad
A_{23},
\qquad
A_{13}.
\]
In quadratic degree, the formal Arnold--Orlik--Solomon relation is
\[
A_{12}A_{23}-A_{12}A_{13}+A_{23}A_{13}=0.
\]
Equivalently, in the normal-form basis
\[
A_{12}A_{23},
\qquad
A_{12}A_{13},
\]
one rewrites
\[
A_{23}A_{13}
=
-A_{12}A_{23}+A_{12}A_{13}.
\]
Substitution gives
\[
A_{12}A_{23}-A_{12}A_{13}
+
\left(-A_{12}A_{23}+A_{12}A_{13}\right)=0.
\]

\begin{theorem}[Lean formalized OS normal form]
The Lean file \texttt{proofs/NonIsoConf3OrlikSolomon.lean} proves that the above Arnold--Orlik--Solomon relation reduces to zero in a rank-two quadratic normal form.
\end{theorem}

The corresponding formal rank bookkeeping is
\[
\dim H^0_{\mathrm{formal}}=1,
\qquad
\dim H^{\mathrm{gen}}_{\mathrm{formal}}=3,
\qquad
\dim H^{\mathrm{quad}}_{\mathrm{formal}}=2.
\]

\begin{theorem}[Lean formalized rank skeleton]
The formal slot-rank function satisfies
\[
\operatorname{rank}(0)=1,
\qquad
\operatorname{rank}(\mathrm{generator})=3,
\qquad
\operatorname{rank}(\mathrm{quadratic})=2,
\qquad
\operatorname{rank}(\mathrm{other})=0.
\]
\end{theorem}

\section{SymPy Audit}
For $D=4$, the SymPy witness uses
\[
q(z)=z_0^2+z_1^2+z_2^2+z_3^2.
\]
For each pair $ij$, it forms
\[
q_{ij}=q(x_i-x_j)
\]
and the logarithmic one-form
\[
d\log q_{ij}=\frac{dq_{ij}}{q_{ij}}.
\]
The script symbolically verifies
\[
d(d\log q_{ij})=0.
\]
It also verifies the Orlik--Solomon normal-form relation and the cooperad generator splitting for the cluster $(12)|3$.

The audit files are:
\begin{itemize}[leftmargin=2em]
\item \texttt{proofs/quadratic\_configuration3.py},
\item \texttt{proofs/non\_iso\_conf3\_orlik\_solomon.py}.
\end{itemize}

\section{Finite-Field Point-Count Fingerprint in Dimension Four}
For the reduced dimension-four model over a finite field, define
\[
U_4(\mathbb F_p)=
\{(a,b)\in \mathbb F_p^4\times \mathbb F_p^4:
q(a)q(b)q(a-b)\neq 0\},
\qquad
q(z)=z_0^2+z_1^2+z_2^2+z_3^2.
\]
A finite-field audit computes this count by a cyclic convolution/FFT method rather than by a general D-module Groebner pipeline.  The observed and Lean-recorded polynomial is
\[
\#U_4(\mathbb F_p)
=
p^2(p-1)^2(p+1)(p^3-2p^2-p+3).
\]
Equivalently,
\[
\#U_4(\mathbb F_p)
=
p^8-3p^7+7p^5-4p^4-4p^3+3p^2.
\]
The Python witness verifies this formula for
\[
p=3,5,7,11,13,17,19,23,29,31.
\]
The Lean file \texttt{proofs/NonIsoConf3QuadricD4PointCount.lean} records the polynomial and proves the sample values
\[
\#U_4(\mathbb F_3)=1296,
\qquad
\#U_4(\mathbb F_5)=175200,
\]
\[
\#U_4(\mathbb F_7)=3400992,
\qquad
\#U_4(\mathbb F_{11})=156961200.
\]

\begin{remark}[Interpretation]
If one proves the required polynomial-count and purity or mixed-Tate comparison theorem, then the count polynomial becomes a candidate compactly supported $E$-polynomial by the substitution $p\mapsto uv$:
\[
E_c(U_4;u,v)
=(uv)^2(uv-1)^2(uv+1)((uv)^3-2(uv)^2-uv+3).
\]
This bridge is not proved in the current Lean kernel.  It is recorded as a typed socket:
\[
\texttt{PointCountToCohomologySocket}.
\]
\end{remark}

\section{Theorem-Honest Boundary}
The current formalization proves the finite skeleton
\[
\boxed{
\text{translation bookkeeping, edge generators, degrees, OS relation, rank slots, and cooperad edge maps.}
}
\]
It does \emph{not} yet prove the analytic de Rham comparison theorem
\[
H^\bullet_{\dR}\!
\left(F_3^{\mathrm{noniso}}(\mathbb C^D,q)\right)
\cong
\mathcal A_{\mathrm{formal}}(K_3,q).
\]
That comparison, together with the construction of the quadric angular class $\beta$, the Brieskorn/Orlik--Solomon comparison theorem for these quadric divisors, and the full cooperad theorem, remains socketed.

\section{Lean Artifacts}
The verified Lean artifacts are:
\begin{itemize}[leftmargin=2em]
\item \texttt{QuadraticConfiguration3.quadratic\_configuration3\_synthesis},
\item \texttt{NonIsoConf3OrlikSolomon.arnold\_orlik\_solomon\_relation},
\item \texttt{NonIsoConf3OrlikSolomon.cohomology\_rank\_synthesis},
\item \texttt{NonIsoConf3OrlikSolomon.non\_iso\_conf3\_os\_cooperad\_synthesis},
\item \texttt{NonIsoConf3QuadricD4PointCount.point\_count\_sample\_synthesis},
\item \texttt{NonIsoConf3QuadricD4EPolynomial.countPolynomialZ\_expand},
\item \texttt{NonIsoConf3QuadricD4EPolynomial.candidateEc\_expand},
\item \texttt{NonIsoConf3QuadricD4EPolynomial.e\_polynomial\_conditional\_synthesis}.
\end{itemize}

The full repository build completed successfully with these files registered as Lean roots.

\section*{Conclusion}
We have formalized the finite algebraic target for the three-point non-isotropic configuration problem over a complex quadric.  The result is not a completed de Rham cohomology theorem; rather, it is a compiler-checked skeleton into which the analytic comparison theorem should land.  In the project terminology: the finite combinatorial spine compiles, while the analytic de Rham and cooperad comparison remains a typed socket.

\end{document}
